\begin{lemma}
\label{lemma-finite-projective}
Let $R$ be a ring and let $M$ be an $R$-module.
The following are equivalent
\begin{enumerate}
\item $M$ is finitely presented and $R$-flat,
\item $M$ is finite projective,
\item $M$ is a direct summand of a finite free $R$-module,
\item $M$ is finitely presented and
for all $\mathfrak p \in \Spec(R)$ the
localization $M_{\mathfrak p}$ is free,
\item $M$ is finitely presented and
for all maximal ideals $\mathfrak m \subset R$ the
localization $M_{\mathfrak m}$ is free,
\item $M$ is finite and locally free,
\item $M$ is finite locally free, and
\item $M$ is finite, for every prime $\mathfrak p$ the module
$M_{\mathfrak p}$ is free, and the function
$$
\rho_M : \Spec(R) \to \mathbf{Z}, \quad
\mathfrak p
\longmapsto
\dim_{\kappa(\mathfrak p)} M \otimes_R \kappa(\mathfrak p)
$$
is locally constant in the Zariski topology.
\end{enumerate}
\end{lemma}

\begin{proof}
First suppose $M$ is finite projective, i.e., (2) holds.
Take a surjection $R^n \to M$ and let $K$ be the kernel.
Since $M$ is projective,
$0 \to K \to R^n \to M \to 0$ splits.
Hence (2) $\Rightarrow$ (3).
The implication (3) $\Rightarrow$ (2) follows from the fact that
a direct summand of a projective is projective, see
Lemma \ref{lemma-characterize-projective}.

\medskip\noindent
Assume (3), so we can write $K \oplus M \cong R^{\oplus n}$.
So $K$ is a  direct summand of $R^n$ and thus finitely generated.
This shows $M = R^{\oplus n}/K$ is finitely presented.
In other words, (3) $\Rightarrow$ (1).

\medskip\noindent
Assume $M$ is finitely presented and flat, i.e., (1) holds.
We will prove that (7) holds. Pick any prime $\mathfrak p$ and
$x_1, \ldots, x_r \in M$ which map to a basis of
$M \otimes_R \kappa(\mathfrak p)$. By
Nakayama's lemma (in the form of Lemma \ref{lemma-NAK-localization})
these elements generate $M_g$ for some $g \in R$, $g \not \in \mathfrak p$.
The corresponding surjection $\varphi : R_g^{\oplus r} \to M_g$
has the following two properties: (a) $\Ker(\varphi)$ is a finite
$R_g$-module (see Lemma \ref{lemma-extension})
and (b) $\Ker(\varphi) \otimes \kappa(\mathfrak p) = 0$
by flatness of $M_g$ over $R_g$ (see
Lemma \ref{lemma-flat-tor-zero}).
Hence by Nakayama's lemma again there exists a $g' \in R_g$ such that
$\Ker(\varphi)_{g'} = 0$. In other words, $M_{gg'}$ is free.

\medskip\noindent
A finite locally free module is a finite module, see
Lemma \ref{lemma-cover},
hence (7) $\Rightarrow$ (6).
It is clear that (6) $\Rightarrow$ (7) and that (7) $\Rightarrow$ (8).

\medskip\noindent
A finite locally free module is a finitely presented module, see
Lemma \ref{lemma-cover},
hence (7) $\Rightarrow$ (4).
Of course (4) implies (5).
Since we may check flatness locally (see
Lemma \ref{lemma-flat-localization})
we conclude that (5) implies (1).
At this point we have
$$
\xymatrix{
(2) \ar@{<=>}[r] & (3) \ar@{=>}[r] & (1) \ar@{=>}[r] &
(7)  \ar@{<=>}[r] \ar@{=>}[rd] \ar@{=>}[d] & (6) \\
& & (5) \ar@{=>}[u] & (4) \ar@{=>}[l] & (8)
}
$$

\medskip\noindent
Suppose that $M$ satisfies (1), (4), (5), (6), and (7).
We will prove that (3) holds. It suffices
to show that $M$ is projective. We have to show that $\Hom_R(M, -)$
is exact. Let $0 \to N'' \to N \to N'\to 0$ be a short exact sequence of
$R$-module. We have to show that
$0 \to \Hom_R(M, N'') \to \Hom_R(M, N) \to
\Hom_R(M, N') \to 0$ is exact.
As $M$ is finite locally free there exist a covering
$\Spec(R) = \bigcup D(f_i)$ such that $M_{f_i}$ is finite free.
By
Lemma \ref{lemma-hom-from-finitely-presented}
we see that
$$
0 \to \Hom_R(M, N'')_{f_i} \to \Hom_R(M, N)_{f_i} \to
\Hom_R(M, N')_{f_i} \to 0
$$
is equal to
$0 \to \Hom_{R_{f_i}}(M_{f_i}, N''_{f_i}) \to
\Hom_{R_{f_i}}(M_{f_i}, N_{f_i}) \to
\Hom_{R_{f_i}}(M_{f_i}, N'_{f_i}) \to 0$
which is exact as $M_{f_i}$ is free and as the localization
$0 \to N''_{f_i} \to N_{f_i} \to N'_{f_i} \to 0$
is exact (as localization is exact). Whence we see that
$0 \to \Hom_R(M, N'') \to \Hom_R(M, N) \to
\Hom_R(M, N') \to 0$ is exact by
Lemma \ref{lemma-cover}.

\medskip\noindent
Finally, assume that (8) holds. Pick a maximal ideal $\mathfrak m \subset R$.
Pick $x_1, \ldots, x_r \in M$ which map to a $\kappa(\mathfrak m)$-basis of
$M \otimes_R \kappa(\mathfrak m) = M/\mathfrak mM$. In particular
$\rho_M(\mathfrak m) = r$. By
Nakayama's Lemma \ref{lemma-NAK}
there exists an $f \in R$, $f \not \in \mathfrak m$ such that
$x_1, \ldots, x_r$ generate $M_f$ over $R_f$. By the assumption that
$\rho_M$ is locally constant there exists a $g \in R$, $g \not \in \mathfrak m$
such that $\rho_M$ is constant equal to $r$ on $D(g)$. We claim that
$$
\Psi : R_{fg}^{\oplus r} \longrightarrow M_{fg}, \quad
(a_1, \ldots, a_r) \longmapsto \sum a_i x_i
$$
is an isomorphism. This claim will show that $M$ is finite locally
free, i.e., that (7) holds. To see the claim
it suffices to show that the induced map on localizations
$\Psi_{\mathfrak p} : R_{\mathfrak p}^{\oplus r} \to M_{\mathfrak p}$
is an isomorphism for all $\mathfrak p \in D(fg)$, see
Lemma \ref{lemma-characterize-zero-local}.
By our choice of $f$ the map $\Psi_{\mathfrak p}$
is surjective. By assumption (8) we have
$M_{\mathfrak p} \cong R_{\mathfrak p}^{\oplus \rho_M(\mathfrak p)}$
and by our choice of $g$ we have $\rho_M(\mathfrak p) = r$.
Hence $\Psi_{\mathfrak p}$ determines a surjection
$R_{\mathfrak p}^{\oplus r} \to
M_{\mathfrak p} \cong R_{\mathfrak p}^{\oplus r}$
whence is an isomorphism by
Lemma \ref{lemma-fun}.
(Of course this last fact follows from a simple matrix argument also.)
\end{proof}
